\section{Salida de OpenAI y filosofía del talento}

\subsection{La verdadera razón para fundar Anthropic}

Dario trabajó en OpenAI cerca de 5 años, y los últimos dos como vicepresidente de investigación. Él e Ilya Sutskever definieron juntos la dirección de investigación en 2016-2017, y Dario impulsó con todo el escalamiento (Scaling): GPT-2, GPT-3, RLHF, debate, amplificación, interpretabilidad.

Sobre la razón de su salida, aclaró dos malentendidos comunes:
\begin{itemize}
    \item \textbf{No} fue por la inversión de Microsoft (esa afirmación es incorrecta)
    \item \textbf{No} fue por la comercialización (él mismo construyó GPT-3 e impulsó su comercialización)
\end{itemize}

La verdadera razón fue tener una visión distinta de «cómo hacerlo»: más cautela, más franqueza, más honestidad, construir confianza. \textit{«¿Cómo puede la seguridad ser algo más que lo que decimos para contratar gente?»}

\begin{importantbox}{«Ve a hacer tu visión»}
\textit{«Discutir tu propia visión dentro de la organización de alguien más es sumamente ineficiente... Lo más productivo es irte, a hacer un experimento limpio.»}

Esta filosofía también se refleja en la estrategia Race to the Top: \textit{«al final no importa quién gane, mientras todos aprendan las buenas prácticas de los demás.»} El verdadero enemigo es el Race to the Bottom.

Anthropic es \textit{«un grupo de personas imperfectas, apuntando de forma imperfecta a un ideal que nunca podrá realizarse a la perfección.»}
\end{importantbox}

\subsection{Densidad de talento contra cantidad de talento}

\textit{«Esta idea es más correcta cada mes que pasa.»}

El experimento mental de Dario: 100 personas sumamente inteligentes y alineadas en dirección $>$ 1000 personas de las cuales 200 son excelentes y 800 son empleados comunes de una gran empresa. La razón:

\begin{itemize}
    \item \textbf{Densidad alta} = confianza, motivación mutua, todos ven el objetivo mayor
    \item \textbf{Densidad baja} = se necesitan procesos, barandales, luchas políticas, cada quien por su lado
\end{itemize}

Anthropic pasó de 300 a 800 personas con rapidez, y luego, de 800 a cerca de 950, se desaceleró de forma notoria. Dario considera que hay un «punto de inflexión» alrededor de las 1,000 personas, a partir del cual hace falta más cautela. Contrató a muchos físicos teóricos porque \textit{«aprenden las cosas muy rápido.»}

Cita a Steve Jobs: \textit{«La gente de tipo A quiere mirar a su alrededor y ver a otra gente de tipo A.»} Ver a personas que no persiguen la misión con todo su esfuerzo \textit{«es desalentador.»}

\subsection{Qué distingue a un buen investigador de IA}

Primera cualidad: \textbf{mente abierta}. La propia experiencia de Dario: vio los mismos datos que todos, no era mejor programador, solo estaba \textit{«dispuesto a ver las cosas con ojos nuevos.»}

\textit{«El descubrimiento de la Scaling Hypothesis fue simple y hasta tonto. Cualquiera pudo haberlo hecho.»} Pero, en los hechos, solo un «puñado de personas» hizo que el campo entero llegara a esta comprensión.

\begin{knowledgebox}{Consejo para los jóvenes}
\begin{itemize}
    \item Empieza directamente a jugar con los modelos: \textit{«estos modelos son algo nuevo que nadie entiende realmente»}
    \item Explora áreas poco estudiadas: interpretabilidad (solo unas 100 personas, contra 10,000 en la dirección de arquitecturas), aprendizaje de largo plazo, evaluación, múltiples agentes
    \item \textit{«Patina hacia donde va a estar el disco»}
    \item La experiencia a veces es una desventaja: la mente abierta suele venir de lo fresca que resulta el área para uno
\end{itemize}
\end{knowledgebox}

\subsection{Resumen de la sección}

La fundación de Anthropic surge de la filosofía de «ir a hacer tu propia visión». La densidad de talento es la ventaja competitiva central: un equipo de alta densidad funciona de manera eficiente sin necesidad de procesos pesados. La cualidad central de un buen investigador de IA es la mente abierta, no la experiencia técnica.
